\section{Finite derivations for the suppositional account}
\label{sec:annex-p19-rationality-6-2}

\begin{annexitem}{annex:p19.r62.material}{thm:p19-r62-material}{Material and conditional-probability assessments}
\begin{proof}
The event \(\neg A\lor B\) is the disjoint union of \(\neg A\)
and \(A\land B\), so its probability is \(1-a+ar\).
Subtracting \(r\) gives \((1-a)(1-r)\). Both factors are
nonnegative, and their product is zero exactly when one factor is zero.
Four equal masses give \(a=1/2\), \(P(A\land B)=1/4\),
and hence the values \(1/2\) and \(3/4\).

For a fair toss, half of the toss-event mass gives heads, so \(r=1/2\)
independently of \(a>0\). For the double-headed toss, no toss-event mass
gives tails, so \(r=0\). Substitution yields \(1-a/2\) and
\(1-a\). These are compatible finite models: choose any distribution
on the no-toss cases to complete mass \(1-a\). The limiting statement
is about a sequence of positive \(a\); it does not assign a ratio to
the endpoint \(a=0\).
\end{proof}
\end{annexitem}

\begin{annexitem}{annex:p19.r62.jeffrey-value}{thm:p19-r62-jeffrey-value}{Expectation and the Jeffrey table}
\begin{proof}
The three disjoint regions defining \(X_r\) give
\[
 E_P(X_r)=P(A\land B)+rP(\neg A).
\]
Thus \(E_P(X_r)=r\) is equivalent to
\(P(A\land B)=rP(A)\). For \(P(A)>0\), division establishes
both directions. If \(P(A)=0\), the subset \(A\land B\) has
mass zero and \(P(\neg A)=1\), making the expectation \(r\)
for every \(r\). The construction is well defined because its
three regions partition the finite space and all its values lie in
\([0,1]\). It supplies no uniqueness in the null case and no
general family-wise betting-coherence result beyond its statement.
\end{proof}
\end{annexitem}

\begin{annexitem}{annex:p19.r62.intervals}{thm:p19-r62-intervals}{Sharp bounds and uncertainty preservation}
\begin{proof}
Order the atoms as \(AB,A\neg B,\neg AB,\neg A\neg B\).
With the first mass \(x\), the mass of \(B\) is \(x+z\),
where \(0\leq z\leq1-x\); every such value is realized by
masses \((x,0,z,1-x-z)\). This proves elimination's interval.
For \(x>0\), the antecedent mass ranges over \([x,1]\),
so \(P(B\mid A)=x/P(A)\) ranges over \([x,1]\).
For a target \(r\) in this interval, choose masses
\((x,x/r-x,0,1-x/r)\). They are nonnegative, sum to one,
and realize the claimed ratio. Consequently one-premise centering
never lowers confidence below its premise.

For given marginals \(a>0,b\), let the first mass be \(z\).
The four masses must then be
\[
 (z,a-z,b-z,1-a-b+z).
\]
All are nonnegative exactly when
\(\max(0,a+b-1)\leq z\leq\min(a,b)\).
Every point in that interval defines a law with the specified marginals.
Dividing by \(a\) establishes both necessity and attainability of the
stated conditional bounds. If \(a+b-1\geq0\), the lower bound
\((a+b-1)/a\) is at least \(a+b-1\) because \(a\leq1\).
If \(a+b-1<0\), any nonnegative conditional assessment is also at
least \(a+b-1\). Hence in both cases
\[
 1-P(B\mid A)\leq(1-a)+(1-b),
\]
which proves two-premise centering's p-validity. Elimination's p-validity
is the already proved inequality \(P(B)\geq x\).

Finally, when \(a,r\) are fixed, the first two masses are \(ar\)
and \(a(1-r)\); splitting the remaining \(1-a\) between the
last two atoms gives exactly \([ar,ar+1-a]\) for \(P(B)\).
For any \(v\) in that interval, masses
\((ar,a(1-r),v-ar,1-a-v+ar)\) realize it. Since
\(ar-(a+r-1)=(1-a)(1-r)\geq0\), all such laws satisfy
\[
 1-P(B)\leq(1-a)+(1-r),
\]
the MP uncertainty budget. Substitution gives the source's numerical
interval and certain-antecedent case. The theorem does not extend its
ratio-based centering claims to \(a=0\) or to \(x=0\) with an
undefined conditional assessment; those require a separately specified
conditional-event framework.
\end{proof}
\end{annexitem}

\begin{annexitem}{annex:p19.r62.update}{thm:p19-r62-update}{Invariant updating and the missing-information counterexample}
\begin{proof}
The prior has positive mass in both cells, so its conditional laws
\(P_1(\cdot\mid F)\) and \(P_1(\cdot\mid\neg F)\) are
well-defined probability laws. A convex mixture with weights \(t\)
and \(1-t\) is nonnegative, normalized and finitely additive.
Its value at \(F\) is \(t\). For \(0<t<1\), dividing its
restrictions to the two cells by their masses recovers the prior
conditional laws. At \(t=1\), the law is \(P_1(\cdot\mid F)\),
so its value at \(H\) is \(r\); the other cell is null.

For fixed \(r,t\), varying \(q\in[0,1]\) in
\(tr+(1-t)q\) yields the stated interval, including both endpoints.
To establish feasibility under just those inputs, choose any prior
cell mass \(s\in(0,1)\) and prior atoms
\((sr,s(1-r),(1-s)q,(1-s)(1-q))\). This realizes every
allowed \(q\). A further specified prior marginal may exclude some
of those laws, which explains the sharpness qualification.

If only the new \(F\) marginal \(t\) is imposed, then for any
\(v\in[0,1]\) the new atoms
\((tv,t(1-v),(1-t)v,(1-t)(1-v))\) form a probability law
with \(P_2(F)=t\) and \(P_2(H)=v\).
Therefore the marginal does not select \(v\) without invariance or
other information constraints. This explicit family proves
underdetermination; it does not merely give the missing update a name.
\end{proof}
\end{annexitem}
